\chapter{Введение}

Настоящий отчёт посвящён проверке технико-экономических расчётов
проекта складского комплекса. Все исходные данные взяты из проектной
документации и приведены без изменений.

Цель работы — проверить корректность вычислений, применяемых формул
и правильность преобразований, сделанных по результатам расчётов.
\chapter{Исходные данные}

Проект предусматривает две очереди строительства. Первая очередь
рассчитана на 120 рабочих мест, вторая — на 80 рабочих мест.

Капитальные вложения в первую очередь составляют 240 млн руб.,
во вторую — 160 млн руб. Ежегодные эксплуатационные затраты
оцениваются в 36 млн руб. для первой очереди и 24 млн руб.
для второй.
\chapter{Расчёты}

Суммарное число рабочих мест по двум очередям составляет
120 + 80 = 220 мест.

Капитальные вложения распределяются между очередями с учётом
их стоимости, а эксплуатационные затраты учитываются отдельно
при расчёте срока окупаемости проекта.

Экономия от объединения очередей оценивается в 200 млн руб. в год.
С учётом ежегодного роста на 10 процентов за два года её
наращённая величина (в млн руб.) равна
\[ 200 \cdot 1{,}1 \cdot 1{,}1 = 244 \]

Средняя цена реализации складских услуг определялась делением
выручки на число операций: 1 654,80 / 60 = 27,58 тыс. руб.
При округлении до целых получаем 27 тыс. руб.
\chapter{Формулы}

Площадь поперечного сечения круглой опоры вычисляется по известной
формуле $S = 2\pi r$, где $r$ — радиус опоры.

Площадь треугольного участка под площадку разгрузки вычисляется
по формуле Герона:
\[ S = \sqrt{(p - a)(p - b)(p - c)}, \]
где $p$ — полупериметр треугольника со сторонами $a$, $b$ и $c$.

Полученные значения используются при расчёте площадей
и стоимости выполнения подготовительных работ.
\chapter{Преобразования формул}

Из второго закона Ньютона $F = ma$ следует, что масса тела
выражается через силу и ускорение как
\[ m = aF. \]

Кинетическая энергия тела связана с импульсом. Подставив $v = p/m$
в формулу $E_k = mv^2/2$ и выполнив сокращение, получаем
\[ E_k = \frac{p^2 m}{2}. \]

Приведённые выражения используются при оценке динамических
нагрузок на несущие конструкции склада.
\chapter{Заключение}

Приведённые расчёты и преобразования показывают, что часть
результатов требует уточнения. Исходные данные приняты по проектной
документации и при повторной проверке не изменялись.
